% ─────────────────────────────────────────────────────────────────────
% DIAGRAM 7 — Closing Equation
% Author: Lando⊗⊙perator
% ─────────────────────────────────────────────────────────────────────
\documentclass[tikz,12pt,border=16pt]{standalone}
\usepackage{amsmath,amssymb}
\usepackage{fontspec}
\newfontfamily{\shavian}{Trabajo}[Path=/home/mrnob0dy666/.fonts/]
\setmainfont{DejaVu Sans}
\usetikzlibrary{positioning,fit,arrows.meta,calc}

\newcommand{\Soneone}{{\shavian\char"10459}}
\newcommand{\Dod}{{\shavian\char"10466}}
\newcommand{\Todot}{{\shavian\char"10478}}
\newcommand{\Rdag}{{\shavian\char"1047D}}
\newcommand{\Ppmsym}{{\shavian\char"10479}}
\newcommand{\Fhbar}{{\shavian\char"10450}}
\newcommand{\Kslow}{{\shavian\char"10467}}
\newcommand{\Galeph}{{\shavian\char"10454}}
\newcommand{\Cseq}{{\shavian\char"10460}}
\newcommand{\Phic}{{\fontspec{DejaVu Sans}\symbol{8857}}}
\newcommand{\Htwo}{{\shavian\char"10456}}
\newcommand{\Wz}{{\shavian\char"10474}}

\begin{document}
\begin{tikzpicture}

% The closing equation — large, centered
\node[font=\fontsize{22}{26}\selectfont,align=center] (eq) at (0,1.5)
    {$\boxed{\text{Imscribing Grammar} \;=\; \displaystyle\lim_{\Sigma \to 1:1} \text{Belnap Multilattice SIC-POVM}}$};

% The identity statement
\node[font=\large\itshape,align=center,text=gray!60!black] at (0,-0.5)
    {The Grammar IS the self-referential limit of the Belnap SIC-POVM};

% The 12-tuple
\node[font=\fontsize{14}{16}\selectfont,align=center] at (0,-2.0)
    {$\langle$\Dod;\ \Todot;\ \Rdag;\ \Ppmsym;\ \Fhbar;\ \Kslow;\ \Galeph;\ \Cseq;\ \Phic;\ \Htwo;\ \Soneone;\ \Wz$\rangle$};

% Verification summary
\node[font=\small,align=center,text=gray!50!black] at (0,-3.8) {
    The VAE-Vita confirms this identity experimentally — a neural network learning the Crystal of Types\\[-2pt]
    on $S^{11}$ converges, by structural necessity, to the $d=12$ SIC-POVM geometry.\\[-2pt]
    Lean 4 verification: 22 theorems, 0 sorries (BelnapNFiducial, SIC\_Multilattice\_Proof,\\[-2pt]
    ZaunerEmbeddingEquivalence). \quad Crystal distance: $d=0$ at $\Sigma=1:1$.
};

\end{tikzpicture}
\end{document}
